\cotefile{84}
\scanpage{2}

\diagram{\begin{tikzcd}[ampersand replacement=\&]
R_1 \arrow[r,shift left,"\deleted{\unsure{\alpha}}\beta"] \& R_0 \arrow[l,shift left,"\deleted{\unsure{\beta}}\alpha"]
\end{tikzcd}}

\deleted{\unsure{hom. de monoïde|hom. de monoïdes}} \quad $R_0,R_1$ monoïdes

\[
\begin{aligned}
\alpha:R_0\longrightarrow R_1&\quad\text{hom. de monoïde}\quad (L,\delta\in\Gamma(L^{\otimes2}),T_0\in\Gamma L_0)\longmapsto (L,\Gamma)\\
\beta:R_1\longrightarrow R_0&\quad\text{\unsure{applicat. ensembl.|application ensem.|multiplicat.} (pas comp. avec produits)}\quad (L,\delta)\longmapsto(L,4\delta,0)
\end{aligned}
\]

\diagram{\begin{tikzcd}[ampersand replacement=\&]
R_\ast\arrow[d] \& {}=R_0\amalg R_1\\
\mathbb Z/2\mathbb Z \&
\end{tikzcd}}

loi de composition \inserted{$(x,y)\mapsto x\wedge y$ comm., ass., unitaire,} comp. avec $R_\ast\longrightarrow\mathbb Z/2\mathbb Z$

\[
\text{i.e.}\quad
\begin{aligned}
R_0\cdot R_0&\subset R_0\\
R_1 R_1&\subset R_0\\
R_1\cdot R_0=R_0\cdot R_1&\subset R_1
\end{aligned}
\]

\[
\begin{aligned}
x,y\in R_0&\Longrightarrow x\wedge y=x\cdot y\\
x,y\in R_1&\Longrightarrow x\wedge y=\beta(x\cdot y)\qquad\text{(pas néc. égal à $\beta(x)\,\beta(y)$)}\\
x\in R_1,\ y\in R_0&\qquad x\wedge y=y\wedge x=\deleted{\unsure{\alpha(x)y}}\,x\alpha(y)
\end{aligned}
\]

\[(x\wedge y)\wedge z=x\wedge(y\wedge z)\]

\[
\begin{aligned}
\text{a)}\quad x,y,z\in R_0&\quad\text{OK}\quad\deleted{\text{(ass. dans $R_0$)}}\qquad\text{ass. dans $R_0$}\\
\text{b)}\quad x\in R_1,\ y,z\in R_0&\quad (x.\alpha(y)).\alpha(z)=x.\underbrace{\alpha(y.z)}_{(\alpha(y)\alpha(z))}\quad\text{mult. de $\alpha$ + ass. dans $R_1$}\\
\text{c)}\quad y\in R_1,\ x,z\in R_0&\quad (\alpha(x)y)\alpha(z)=\alpha(x)(y\alpha(z))\quad\text{ass. dans $R_1$}\\
\text{d)}\quad z\in R_1,\ x,y\in R_0&\quad \alpha(x.y).z=\alpha(x)(\alpha(y).z)\quad\text{mult. de $\alpha$ + ass. de $R_1$}
\end{aligned}
\]

\[
\begin{aligned}
\text{b$'$)}\quad x\in R_0,\ y,z\in R_1&\quad\beta((\alpha(x).y).z)\overset?=x.\beta(y.z)\quad (\ast)\\
\text{c$'$)}\quad y\in R_0,\ x,z\in R_1&\quad\beta((x\alpha(y)).z)=\beta(x.(\alpha(y).z))\quad\text{ass. dans $R_1$}\\
\text{d$'$)}\quad z\in R_0,\ x,y\in R_1&\quad\underset{\Vert}{\beta(x.y)z}\overset?=\underset{\Vert}{\beta(x(y\alpha(z)))}\quad\text{comme $(\ast)$ + comm. et ass. dans $R_1$}\\
&\quad z\,\beta(y.x)\qquad\beta((\alpha(z)y)x)\\
\text{a$'$)}\quad x,y,z\in R_1&\quad\alpha(\beta(x.y)).z\overset?=x\,\alpha(\beta(y.z))
\end{aligned}
\]

Donc il faut vérifier

\[
\beta(\alpha(x).\underbrace{yz})=x\,\beta(\underbrace{yz}_{u})\qquad\text{i.e.}\tag*{\textcircled{1}}
\]

\margin{(faisant $y$ ou $z=1$)\quad $\boxed{\beta(\alpha(x)u)=x\beta(u)}$\quad $\swarrow$ $\boxed{\beta(\alpha(x))=x.\beta(1)}$}

\[
\begin{aligned}
(\alpha.\beta(xy))z&=x\,\alpha\beta(yz)\\
(\alpha\beta(x))z&=x\,(\alpha\beta(z))
\end{aligned}\tag*{\textcircled{2}}
\]

\margin{si $y=1$ cela donne et si $z=1$\quad $\boxed{\alpha\beta(x)=x.\underbrace{\alpha\beta(1)}_{\unsure{\xi|\eta}}}$}

Et \textcircled{2} devient alors $(xy\zeta)z=x(yz\zeta)$ OK par ass. et comm. de $R_1$

\diagram{\begin{tikzcd}[ampersand replacement=\&,column sep=small]
\boxed{\begin{aligned}\text{\textcircled{A}}&\quad\beta(\alpha(x)u)=x\beta(u)\\
\text{\textcircled{B}}&\quad\alpha\beta(x)=x.\underbrace{\eta}_{(\underline{\mathcal O}_X,4.\underline{1}_{\mathcal O_X})}\ \text{où }\eta=\alpha\beta(1)=\alpha(\zeta)\quad\zeta=\underbrace{\beta(1)}_{(\underline{\mathcal O}_X,4\underline1_{\mathcal O_X},0_{X_0})}\end{aligned}}\arrow[d,Rightarrow]\arrow[dr,Rightarrow] \&\\
\boxed{\beta\alpha(x)=x\zeta} \& \boxed{(\beta x)(\beta y)=\zeta\,\beta(x.y)}
\end{tikzcd}}

OK

\[
\begin{aligned}
&\deleted{\beta(1)\beta(x.y)1=\beta}\\[-2pt]
&\boxed{\underbrace{\beta(x)\beta(y)}_{\Vert\,\text{\textcircled{A}}}\overset?=\deleted{\unsure{\beta(xy)}}\underbrace{\zeta}_{\beta(1)}\cdot\beta(\overbrace{x.y}^{u})}\\
&\beta(\underbrace{\alpha\beta(x)}_{x\eta\ \text{\textcircled{B}}}y)\qquad\quad\Vert\,\text{\textcircled{A}}\\
&\beta(xy\eta)\qquad
\beta(\underbrace{\underbrace{\alpha(\zeta)}_{\eta}xy}_{\alpha\beta(xy)\ \text{\textcircled{B}}})=\beta\underbrace{\alpha\beta(xy)}_{\beta\unsure{(\beta|1\,\beta)}(xy)\eta)}
\end{aligned}
\]

Vér. \unsure{les cas|sur cas} d'espèces

\[
\begin{aligned}
\underline{x=(L,\delta)}&\qquad\beta(x)=(L,4\delta,0)\qquad\alpha\beta(x)=(L,4\delta)\\
x=(L,\delta,T_0),\ u&=(L',\delta')\qquad\qquad\quad=(L,\delta)(\underline{\mathcal O},4\underline1_X)\\
\alpha(x)&=(L,\delta)\\
\alpha(x).u&=(LL',\delta\delta')\\
\beta(\alpha(x).u)&=(LL',4\delta\delta',0)\\
\beta(u)&=(L',4\delta',0)\qquad x\beta(u)=(LL',4\delta\delta',0)
\end{aligned}
\]

$R_\ast=R\times_{\mathbb Z}\mathbb Z/2\mathbb Z$\qquad On définit

\[
\begin{aligned}
&R_\ast\xrightarrow{\deleted{\unsure{w_n}}} R_0\quad\text{par}\\
&\text{a) si }n=2\nu\qquad R_{2\nu}\xrightarrow{w_n} R_0\qquad w_n(x)=x\zeta^\nu\\
&\hspace{45pt}\deleted{\unsure{R_1}}\quad R_0\\
&\text{b) si }n=2\nu+1\qquad\deleted{n+1=2\nu+2}\qquad w_n(\unsure{x})=\beta(x)\zeta^\nu\\
&\hspace{48pt}\underset{R_1}{\underset{\Vert}{R_n}}
\end{aligned}
\]

On trouve ainsi un \uline{hom} $w:R_\ast\longrightarrow R_0$. \unsure{d'où|Un}
\[w'_1:R_\ast\longrightarrow\mathbb Z\times R_0\]

\[
\begin{aligned}
E&=\bigoplus_{1\leq i\leq n}(L_i,\delta_i)\\
w_1(E)&=\bigwedge_{1\leq i\leq n}(L_i,4\delta_i,0)
=\left(\bigotimes\unsure{\vphantom{L}\zeta} L_i,\underbrace{4^n}_{=(2^n)^2}\underbrace{\textstyle\prod\delta_i}_{\delta},0\right)
\end{aligned}
\]

\[
\begin{aligned}
\operatorname{Arf}(E)&=\left(\bigotimes L_i,\underbrace{2^\nu}_{4^{\nu'}=(2^{\nu'})^2}\underbrace{\textstyle\prod\delta_i}_{\delta},0\right)
\quad\begin{aligned}\nu&=n\ \text{si $n$ pair}\\\nu&=n+1\ \text{si $n$ impair}\end{aligned}\\
\operatorname{Arf}(E)&=w_1(E).\zeta_n\qquad \zeta_n=(\underline{\mathcal O}_X,2^n\mathbin{\unsure{:|.}}1,1
\end{aligned}
\]

\margin{de ttes façons $\nu$ est pair\quad $\nu=2\nu'$\quad $\nu'=\left[\dfrac{n+1}{2}\right]$}

\[
\begin{aligned}
w_1(\unsure{L|L_n})&=(L,4\delta,0)\\
\operatorname{Arf}(L)&=(L,4\delta,0)
\end{aligned}
\]

\margin{\deleted{On peut prendre} \inserted{$\unsure{n}$} NB si $2$ inv. \inserted{sur $X$} alors \deleted{$\unsure{d}\sim\quad\otimes$} \deleted{$\operatorname{Arf}(E)\sim$}}

\[
n-\frac{n+1}{2}=\frac{\unsure{n-1}}2\qquad
\nu''=\begin{cases}\frac n2&\text{si $n$ pair}\\ \frac{n-1}2&\text{si $n$ impair}\end{cases}
\]

\unsure{On a.}

\[
\boxed{c_1(E)=\operatorname{Arf}(E).\zeta^{(n-\nu')=\nu''}}
\qquad
\begin{aligned}
\text{où}\quad\zeta&=(\underline{\mathcal O}_X,4,0)\\
&=\operatorname{Arf}(\mu_{2X}/X)
\end{aligned}
\]

donc $\operatorname{Arf}(E)$ est un inv. plus fin que $c_1(E)$ (si $2$ \uline{non} \uline{inv.} sur $X$)

\[
\begin{aligned}
c_1(E)&\simeq\left(\bigotimes L_i,\unsure{\delta},0\right)\\
\operatorname{Arf}(E)&\simeq\left(\bigotimes L_i,\delta,0\right)
\end{aligned}
\]

Mais $c_1(E)$
% archive sheet number: 1
\scanpage{3}

% printed form
P

Je vais \unsure{ouvrir}

\uline{\unsure{Pourny|Pourry}} \uline{Mr. Pascal}\\

\[
\begin{aligned}
&R_1\qquad R_0\\
R_1\times R_1&\longrightarrow R_0\\
R_0\times R_0&\longrightarrow R_0\\
R_0\times R_1&\longrightarrow R_1\\
R_1\times R_0&\mathrel{\unsure{-|\to}} R_1
\end{aligned}
\]

% printed form
% margin note written sideways
OFF\quad Offsetpapier (hochweiß, holzfrei)\\
Offset paper (super white, woodfree)\\
Papier offset VII (extra blanc, sans bois)\\
Offsetpapier (wit, houtvrij)\quad Made in W.-Germany\\
500\quad Blatt\quad sheets\quad feuilles\quad vel\qquad 70\quad g/m²\quad gsm\quad grs/m²\quad grs/m²\qquad \uline{21} x 29,7\quad cm\quad cms\quad cm\quad cm

$R=R_1\amalg R_0$ est un monoïde unitaire

\deleted{\unsure{et}} $R\xrightarrow{\varepsilon}\mathbb Z/2\mathbb Z$ est un hom de Monoïdes

\[
\begin{aligned}
R_1&=\varepsilon^{-1}(1)\\
R_0&=\varepsilon^{-1}(0)\\
R_{00}&\subset R_0\quad\longrightarrow\text{partie étale}=H^1(X,\mathbb Z/2\mathbb Z)
\end{aligned}
\]

\diagram{\begin{tikzcd}[ampersand replacement=\&,column sep=small]
\&\text{\uline{Impair}}\&(L,\delta,T_0)\\
(L,\delta)\arrow[r,mapsto]\&( \pi\ \unsure{w|\omega})\arrow[u]\arrow[d]\&(L,\delta,T_0)\arrow[r,mapsto]\&(L,\delta)\\
(L,4\delta,0)\&\text{\uline{Pair}}\&(L,\delta)
\end{tikzcd}}

\[\deleted{\unsure{\delta}}\qquad(\underline L,\delta)\qquad\underline L,\]

$R_1\longrightarrow R_0$\qquad ce n'est \uline{pas} un hom de monoïdes\,!

\[
(L,\underset{\substack{\downarrow\\\text{\unsure{discr.}}\\\text{\unsure{divisé}}}}{\delta})\longrightarrow (L,4\delta,0)
\]

\[
\begin{aligned}
&R_1\longrightarrow H^1(\underline L)\qquad \deleted{R_1\times R_0}\\
&\Theta=2\delta\qquad\deleted{2\delta}\qquad\deleted{2\delta,\ 2\delta\quad 2\delta}\\
&\Theta'=2\delta'\qquad\Theta\Theta'=4\delta\delta',\ \deleted{2\delta\delta'}\,0\\
&\Theta\delta'=2\qquad\Theta=\underline{2\delta}.
\end{aligned}
\]

\[
\boxed{\begin{aligned}
(L,\delta)(L',\delta')&=(LL',4\delta\delta',0)\\
(L,\delta)(L',\delta',T'_0)&=(LL',\deleted{\unsure{2}}\delta\delta')\\
(L,\delta,T_0)(L',\delta')&=(LL',\delta\delta')\\
(L,\delta,T_0)(L',\delta',T'_0)&=(LL',\delta\delta',T_0T'_0)
\end{aligned}}
\]

\[
\begin{aligned}
&(LL'L'',4\delta\delta'\delta'',0\ \text{ou }\unsure{T_0T'_0T''_0})\quad \text{si deux ou trois parités impaires}\\
&\text{sinon}\quad LL'L'',\delta\delta'\delta'',\ T_0T'_0T''_0
\end{aligned}
\]

Merci pour bonnes choses \inserted{(\uline{\unsure{conversations|conversation|Commutation|Commissions}})} hier et aujourd hui, et pour aide et assistance\,! (\uline{Avant}-\uline{hier})

Comment on compte pour \uline{figues}
\[13\,\mathrm{kg}=\text{----}\]
Je vais payer les figues et tu me rembourses ce que tu en prends. Comme ça je n'ai pas besoin de peser.

Et \unsure{apportes-moi|apporte-moi} stp mon \uuline{ail}, il va se gâter chez toi\,! \uuline{Demain\,?} \uline{\unsure{Plant}} (\inserted{beaucoup de} manquants un peu partout) et \uline{radis}\,?? (\uline{semer})
% archive sheet number: 2
\scanpage{4}

\[
\begin{array}{ccc}
A_1&A_2&A_3\\
1\ u_1&1\ u_2&1\ u_3
\end{array}
\qquad u_i^2+c_i=0
\]

extensions quadratiques, involutions canoniques $\sigma_1,\sigma_2,\sigma_3$ ($\sigma_i(u_i)=-u_i$) se prolongent à $A$ de façon canonique par $\sigma_1\otimes\mathrm{id}\otimes\mathrm{id}$ etc.

\[
\begin{aligned}
&A_1\otimes A_2\otimes A_3\overset{\unsure{\mathrm{déf}|\mathrm{dfn}|\mathrm{df}}}=A\\
&1\quad u_1\quad u_2\quad u_3\quad u_2u_3\quad u_3u_1\quad u_1u_2\quad u_1u_2u_3
\end{aligned}
\]

% margin note written sideways
\margin{correspondent aux sept hyperplans de $\mathbb F_2^3$, i.e. par \unsure{pol.} \unsure{orth.} aux sept droites de $\mathbb F_2^3$, i.e. aux \deleted{sept éléments non nuls}}

\diagram{$\begin{array}{llll}
k&1&\sigma_1,\sigma_2,\sigma_3,\sigma_2\sigma_3,\sigma_3\sigma_1,\sigma_1\sigma_2,\sigma_1\sigma_2\sigma_3=\tau&\\
A_1&1,u_1&\sigma_2,\sigma_3,\sigma_2\sigma_3&\\
A_2&1,u_2&\sigma_3,\sigma_1,\sigma_3\sigma_1&\\
A_3&1,u_3&\sigma_1,\sigma_2,\sigma_1\sigma_2&\\
B_1&1,u_2u_3&\sigma_1,\sigma_2\sigma_3,\sigma_1\sigma_2\sigma_3=\tau&B_1\simeq A_2\wedge A_3\\
B_2&1,u_3u_1&\sigma_2,\sigma_3\sigma_1,\sigma_2\sigma_3\sigma_{\unsure{1|3}}=\tau&B_2\simeq A_3\wedge A_1\\
B_3&1,u_1u_2&\sigma_3,\sigma_1\sigma_2,\sigma_3\sigma_1\sigma_2=\tau&B_3\simeq A_1\wedge A_2\\
B_0&1,u_1u_2u_3&\sigma_2\sigma_3,\sigma_3\sigma_1,\sigma_1\sigma_2&B_0\simeq A_1\wedge A_2\wedge A_3
\end{array}$}

\diagram{$\left.\begin{array}{lll}
A'_1=(A_{\unsure{2|3}},A_3,B_1)&1,u_2,u_3,u_2u_3&\sigma_1\\
A'_2=(A_3,A_1,B_2)&1,u_3,u_1,u_3u_1&\sigma_2\\
A'_3=(A_1,A_2,B_3)&1,u_1,u_2,u_1u_2&\sigma_3\\
B'_1=(A_1,B_1,B_0)&1,u_1,u_2u_3,u_1u_2u_3=u&\sigma_2\sigma_3\\
B'_2=(A_2,B_2,B_0)&1,u_2,u_3u_1,u_2u_3u_1=u&\sigma_3\sigma_1\\
B'_3=(A_3,B_3,B_{0})&1,u_3,u_1u_2,u_3u_1u_2=u&\sigma_1\sigma_2\\
B'_0=(B_1,B_2,B_3)&1,u_2u_3,u_3u_1,u_1u_2&\sigma_1\sigma_2\sigma_3
\end{array}\right\}$}

correspondent aux sept droites vect. de $(\mathbb F_2^3)^{\deleted{\unsure{\ast}}}$ i.e. aux \deleted{sept} sept éléments non nuls

\[
\begin{aligned}
&A_i\otimes A'_i\xrightarrow{\sim} A\\
&[B_i\subset B'_i]\\
&B_0\otimes B'_0\xrightarrow{\sim} A\quad\text{si les $c_i$ inv. i.e. les $A_i$ étales}
\end{aligned}
\]

\[
\begin{aligned}
A_1&\subset A'_2,A'_3,B'_1\\
A_2&\subset A'_3,A'_1,B'_2\\
A_3&\subset A'_1,A'_2,B'_3\\
B_1&\subset A'_1,B'_1,B'_0\\
B_2&\subset A'_{\unsure{1|2}},B'_{\unsure{1|2}},B'_0\\
B_3&\subset A'_3,B'_3,B'_0\\
B_0&\subset \underset{B_0(A_1)}{\underset{\Vert}{B'_1}},\underset{B_0(A_2)}{\underset{\Vert}{B'_2}},\underset{B_0(A_3)}{\underset{\Vert}{B'_3}}
\end{aligned}
\]

il y a $7+7=14$ situations ``\unsure{bi}quadratiques'' typiques (quand les $A_i$ sont étales\ldots)
% archive sheet number: 3
\scanpage{5}

\mspage{449}

qu'il puisse avoir appris quelque chose de moi. Pourtant, chose remarquable, mon nom est prononcé trois fois dans cet éloge (nullement funèbre pour le coup) de Deligne(``troisième médaille Fields de L'IHES''). Et même dans une périphrase il est fait allusion, avec le vague de rigueur qui doit entourer chaque apparition de ma modeste personne, au fait que j'aurais ``construit la théorie de cohomologi\unsure{e} en géométrie sur un corps quelconque'' - et sûrement encore ``avec la plus grande généralité naturelle'', ça sent la grothendieckerie à plein nez (*). La citation complète du contexte vaut la peine d'être donnée, c'est un petit chef-d'oeuvre du genre:

``Partant de là \handwritten{[}théorie de Hodge classique\handwritten{]} et d'analogies \unsure{l|!}-adiques suggérées par Grothendieck \handwritten{[}on se demande où Gr. a trouvé le temps d'apprendre des choses aussi sérieuses, tout en rédigeant ses vingt volumes de plus grandes généralités naturelles\handwritten{]}, il \handwritten{[}Deligne\handwritten{]} a dégagé la notion de structure de Hodge mixte et en a muni la cohomologie de toute variété algébrique complexe. En cohomologie l-adique, donc \handwritten{[}?\handwritten{]} pour des variétés sur un corps fini, il a prouvé les conjectures de Weil, d'une difficulté proverbiale. Ce résultat a paru d'autant plus surprenant \handwritten{[}!!\handwritten{]} que Grothendieck, après avoir construit la théorie de cohomologie en géométrie sur un corps quelconque \handwritten{[}on se demande bien ce qu'il est e\unsure{n}core allé chercher là\handwritten{]}, avait ramené la conjecture restante \handwritten{[}???\handwritten{]} à une série de conjectures qui sont aujourd'hui aussi inabordables qu'alors.''

En clair, bien loin d'avoir contribué en quoi que ce soit à prouver ce surprenant résultat d'une difficulté si proverbiale, ces grothendieckeries-là (au nom à faire fuir le généraliste-naturaliste le plus endurci) ont été tout juste bonnes à nous encombrer encore de \mbox{c o n j e c t u r e s} comme de juste (il n'en fait jamais d'autres!) et inabordables ce qui plus est (on s'en serait douté), tout autant aujourd'hui que quand il a eu l'idée saugrenue de les faire.

\noindent\rule{.22\linewidth}{.4pt}

(*) (18 mai) Dans l'Eloge Funèbre, il est question de la ``grande attention'' que je portais à la terminologie. Dans l'utilisation d'expressions saugrenues comme ``la plus grande généralité naturelle'' ou ``la théorie de cohomologie en géométri\unsure{e} sur un corps quelconque'', je perçois clairement l'intention de tourner en dérision cette attention.

Le soin extrême que j'accorde aux noms donnés aux choses découle naturell\unsure{e}ment du respect que j'ai pour ces choses, dont le nom est censé exprimer l'esse\unsure{n}ce, ou du moins quelque aspect essentiel. Par les echos qui me parviennent, j'a\unsure{i} été choqué plus d'une fois par l'affectation de dédain qui aujourd'hui semble d\unsure{e} mise vis-à-vis de cette attitude de respect, dédain qui s'exprime parfois par l'usage de noms abracadabrants pour des notions importantes. Voir aussi à ce sujet la note ``La Perversité'' (n°76).
% archive sheet number: 4
\scanpage{6}

\[
\begin{aligned}
1,u&\qquad -u+\lambda\\
1,u&\longrightarrow \bar u+\alpha
\end{aligned}
\]

\[
\begin{aligned}
N(\bar u+\alpha)&=N(\bar u)\quad(=N(u))\qquad
\operatorname{Tr}(\bar u+\alpha)=\operatorname{Tr}\bar u\ (= \operatorname{Tr}u)\\
\Vert&\\
N(\bar u)+\alpha\operatorname{Tr}(\bar u)+\alpha^2&=N(\bar u)\\
\left\{\begin{aligned}\alpha(\underbrace{\operatorname{Tr}\bar u}_{b}+\alpha)&=0\\2\alpha&=0\end{aligned}\right.&\quad\Longrightarrow\alpha=0\quad\text{si $2$ inv.}
\end{aligned}
\]

\[
\begin{aligned}
&\deleted{\lambda+\mu u}\qquad\qquad u\bar u=N(u)\\
&N(v)=1\quad\Longrightarrow\quad v=u\bar u^{-1}\ \deleted{\unsure{\in}}\ (=u^2N(u)^{-1})\\
&\hspace{50pt}u^2+ub +c=0\qquad \operatorname{Tr}u=-b
\end{aligned}
\]

\[
\begin{aligned}
v&=\alpha+\beta u& N(v)&=\alpha^2\deleted{\unsure{+}}-\alpha\beta b+\beta^2c=1\\
\bar v&=(\alpha-\beta b)-\beta u&&=\alpha(\alpha-\beta b)\quad\text{si }c=0\\
v\bar v&=\alpha(\alpha-\beta b)-\beta^{\unsure{2|+}}c\,-\\
v&=\alpha\unsure{-|+}\beta u&v\bar v&=\alpha^2+\beta^2+\alpha\beta\underbrace{(u+\bar u)}_{\operatorname{Tr}u}\\
\bar v&=\alpha+\beta\bar u
\end{aligned}
\]

\[
\begin{aligned}
w&=\lambda1+\mu u\quad\text{si }\deleted{\unsure{N(u)=|N(u)=0}}\,N(u)=0,\ \text{donc }N(w)\ \text{inv.}\\
N(w)&=\lambda^2\deleted{\unsure{+}}-\lambda\mu b=\lambda(\lambda-\mu b)\qquad\Longrightarrow\lambda\ \text{inv.}
\end{aligned}
\]

\[\alpha(\alpha-\beta b)=1\]

\diagram{\begin{tikzcd}[ampersand replacement=\&,column sep=small]
1\arrow[r]\& SW/\mu_{\unsure{2}}\arrow[r,"\sim"]\&W^\ast/G_m\arrow[r]\&1\&\\
\&\deleted{\unsure{w\bar w}}\&\deleted{\unsure{N}}\&\&\\
1\arrow[r]\& SW\arrow[uu]\arrow[r,dash]\& W^\ast\arrow[uu]\arrow[r,"N"]\& G_m\arrow[uu]\arrow[r]\&1\\
1\arrow[r]\&\mu_2\arrow[u]\arrow[r,dash]\&G_m\arrow[u,"{\text{cart.}\atop\text{\unsure{et commut.|évident}}}"']\arrow[r,"{\unsure{\zeta^2|z^2|(\,)^{2}}}"]\&G_m\arrow[u,equal]\arrow[r]\&1\\
\&{}\arrow[u]\&{}\arrow[u]\&{}\arrow[u]\&
\end{tikzcd}}
% archive sheet number: 5
\scanpage{7}

\mspage{444} \handwritten{\unsure{h|4}}

d'avance (qu'il qualifiera de ``vraie'', ou de ``vérité'') n'aura aucune peine à la trouver, et même à la prouver à son entière satisfaction - et sûrement il trouvera bien en chemin tel ou tel autre, si ce n'est toute une foule, tout content de conclure alliance avec lui et de partager convictions et satisfaction. Il est comme le chasseur de papillons, qui part avec dans son filet un beau papillon (empaillé si ça se trouve), et qui le sort tout content ( et à son entière satisfaction) en revenant de sa ``chasse''.

Et il y a celui aussi qui se trouve placé devant un inconnu, comme un enfant nu devant la mer. Quand l'enfant désire la connaître, il entre et la connaît - qu'elle soit tiède ou fraîche, calme ou agitée. Celui qu'attire telle chose inconnue, et qui part pour la connaître, sûrement la connaîtra peu ou prou\unsure{.} Avec ou sans filet, il trouvera le vrai, ou en tous cas. \mbox{d u}\quad vrai. Ses erreur\unsure{s} comme ses trouvailles sont autant d'étapes dans son cheminement, ou pour mieux dire, dans \mbox{s e s}\quad \mbox{a m o u r s} avec ce qu'il désire connaître.

Je sais bien de quoi je parle, car dans ma vie j'ai été abondamment tour à tour, et ce chasseur de papillons, et cet enfant nu. Il n'y a aucune difficult\unsure{é} à distinguer l'un de l'autre. Je doute que les ``critères objectifs'' soient ici d'un grand secours, c'est beaucoup plus simple que çà! Il n'y a qu'à se servir de ses yeux\ldots

Et il n'y a aucune difficulté non plus à distinguer les étapes succéssives\unsure{,} les stades de décantation successifs, dans ce cheminement dont je viens de parle\unsure{r} à partir de cette étape ``morte'' ou nul pressenti \deleted{\unsure{e}}\handwritten{a}ffleurant à la conscience ne fait encore soupçonner ``quelque chose'', au delà d'une certaine surface plate et amorphe que nous présentent des yeux somnolents, et qui à travers des ``éveils'' successifs nous amène vers une appréhension de plus en plus délicate, plus intime, plus complète de ce ``quelque chose''. Il n'est pas de nature essentiellemen\unsure{t} différente, qu'il s'agisse du cheminement dans la découverte des choses mathématiques, ou dans celle de soi et d'autrui. Le sentiment d'une \mbox{p r o g r e s s i o n} dans une \mbox{c o n n a i s s a n c e}, qui \deleted{\unsure{a}}\handwritten{s}'approfondit peu à peu (fût-ce à travers une accumulation d'erreurs, patiemment, inlassablement corrigées) - ce sentiment est aussi irrécusable dans ce dernier cas comme dans l'autre.

Cette \mbox{a s s u r a n c e} - là est l'une des faces d'une disposition intérieure, dont l'autre face est une \mbox{o u v e r t u r e}\quad \mbox{a u}\quad \mbox{d o u t e}: une attitude de curiosité excluant toute crainte, vis à vis de ses propres erreurs, qui permet de les dépister et de les corriger constamment. La condition essentie\unsure{l}le de cette double assise, de cette \mbox{f o i} indispensable pour accueillir le doute comme pour découvrir, est l'absence de toute peur (qu'elle soit apparente
% archive sheet number: 6
\scanpage{9}

Formes \uline{quadratiques binaires et} revêtements quadratiques

1. Forme quadratique à scalaire près, structure de similitude.

1.1. $(X,\underline{\mathcal O}_X)$ désigne un schéma, plus généralement un topos loc. annelé. Soit $E$ un Module sur $X$,

\[\underline{\mathrm{Quad}}(E)\simeq\underline{\mathrm{Hom}}(\Gamma^2(E),E)\]

le faisceau des formes quadratiques sur $E$. Si $E$ est loc. libre de type fini, on a un iso. can.

\[\underline{\mathrm{Quad}}(E)\simeq\underline{\mathrm{Sym}}^2(E^\vee),\]

(en général on a seulement un homomorphisme can.

\[\underline{\mathrm{Sym}}^2(E^\vee)\longrightarrow\underline{\mathrm{Quad}}(E),\]

qui transforme $f.g$ ($f,g\in\Gamma E^\vee$) dans $Q_{\unsure{fg|f,g}}:x\longmapsto f(x)g(x)$ sur $E$ \ldots).

Le faisceau \deleted{$\unsure{\underline{\mathcal O}}$} en groupes $\underline{\mathcal O}_X^\ast$ opère sur $\underline{\mathrm{Quad}}(E)$, d'où un faisceau quotient

\[\underline{\mathrm{Sim}}(E)\simeq\underline{\mathrm{Quad}}(E)/\underline{\mathcal O}^{\ast}.\]

Une section du faisceau quotient s'appelle une \uline{structure de similitude} sur $E$. Une telle section correspond à un sous-faisceau ensembliste de $\underline{\mathrm{Quad}}(E)$, qui \deleted{loc} soit à fibres non vides et sur lequel $\mathcal O^\ast$ opère transitivement. Localement, une telle structure est définie par une section $Q$ de $\underline{\mathrm{Quad}}(E)$ i.e. par une forme quadratique sur $E$, \deleted{\unsure{deux}} deux telles sections $Q$, $Q'$ définissent la $=$ structure de similitude \unsure{sss|ssi} $\exists$ localement $a\in\Gamma\underline{\mathcal O}_X^\ast$, avec $Q'=aQ$.

Le cas le plus intéressant est celui où \deleted{$a\mapsto aQ$}

\[a\longrightarrow a.Q:\underline{\mathcal O}_X\ (\text{sans $\ast$\,!})\longrightarrow\underline{\mathrm{Quad}}(E)\]

est injectif. Dans ce cas, \deleted{si} $Q'$ \deleted{\unsure{est connu de|définit la même}} \inserted{définit la} $=$
% archive sheet number: 8
\scanpage{10}

structure de similitude \unsure{sss|ssi} $\exists\,a\in\Gamma\underline{\mathcal O}_{\unsure{\alpha|X}}^\ast$ (\uline{globalement}) telle que $Q'=aQ$, et cette $a$ est alors unique (\unsure{en quoi résume|en quoi réside|ce qui assure} l'unicité). \deleted{On appelle des} \inserted{Structure de similitude} fidèle : définie loc. par une $Q$ telle que $a\longmapsto aQ$ injectif (\deleted{cela ne dépend pas du choix de} du choix de \deleted{$Q$}, seulement de la structure de similitude).

1.2. \deleted{Soit}\ \deleted{\unsure{par|pour}} $P\subset\underline{\mathrm{Quad}}(E)$ est une structure de similitude, on peut regarder le sous-Module engendré, soit

\[\Delta=\Delta_P\subset\underline{\mathrm{Quad}}(E),\]

et on voit que

\[P\mathrel{\unsure{-|\longrightarrow|\longmapsto}}\Delta_P\]

donne une bijection entre l'ens. des structures de similitudes sur $E$, et l'ens des sous-Modules de $\underline{\mathrm{Quad}}(E)$ qui sont localement monogènes. Les \deleted{\unsure{formes q|formes}} structures \deleted{fidèl} de similitude \uline{fidèles} correspondent aux sous modules

\[\Delta\lhook\joinrel\longrightarrow\underline{\mathrm{Quad}}(E)\]

qui sont \uline{inversibles}.

\margin{\uline{NB} Dans le cas où on a un hom \uline{quelc.} $\Delta\longrightarrow\underline{\mathrm{Quad}}(E)$ (donc non \unsure{forc.} injectif) \unsure{s'in.}\ \deleted{\unsure{est}} \deleted{\unsure{fid}} définissant \deleted{une} \deleted{structure}}

1.3. Si $E$ est \inserted{\unsure{loc\textsuperscript{t}.|loc.}} \unsure{lib.|L.} de type fini, il en \deleted{\unsure{résulte}} est de $=$ de $\underline{\mathrm{Quad}}(E)\simeq\underline{\mathrm{Sym}}^2(E)$. Une structure de similitudes est dite \uline{strictement fidèle} (ou universellement fidèle) si l'hom. précédent est ``universellement injectif'', ou encore (ce qui revient au $=$) injectif sur chaque fibre résiduelle,
% archive sheet number: 9
\scanpage{11}

Les structures de similitude strictement fidèles correspondent donc aux \uline{sous-fibrés vectoriels de rang 1} du fibré Vectoriel $\unsure{W|V|\mathbb W}(\underline{\mathrm{Quad}}(E))$, ou encore aux sections du fibré en espaces projectifs $P^\vee(\underline{\mathrm{Quad}}(E))$.

1.4. Soient $\Delta$ un Module inv. sur $X$, et

\[\deleted{\unsure{\Delta}}\qquad\Delta\xrightarrow{\deleted{\unsure{\alpha}}i}\underline{\mathrm{Quad}}(E)\]

un hom (pas néc. injectif), définissant \inserted{donc} une structure de similitude $\Sigma$ sur $E$. Je dis qu'on a une forme quadratique \inserted{$Q_\Sigma$ ou $Q$} canonique à valeurs dans le module inversible $\Delta^\vee$. Pour la définir, je vais d'abord définir (indép. de la donnée de $i$ ou de $\Sigma$) une forme quadratique canonique

\[Q=Q_E:E\longrightarrow\underline{\mathrm{Quad}}(E)^\vee:\]

\deleted{elle associe} sa valeur, pour une section locale $x$ de $E$, est la forme linéaire

\[Q(x)=\{f\longmapsto f(x)\}\quad\inserted{\text{i.e. }f(x)=\langle f,Q(x)\rangle}\]

\inserted{(pour toute section (locale) $f$ de $\underline{\mathrm{Quad}}(E)$, $x$ de $E$,} sur $\underline{\mathrm{Quad}}(E)$. Celle-ci dépend de $x$ de façon quadratique, i.e. la \deleted{fonction} $E\times E\longrightarrow\underline{\mathrm{Quad}}(E)^\vee$,

\[
\unsure{\text{qui transforme}}\quad x,y\ \text{en}\quad Q(x+y)-Q(x)-Q(y)=\{f\longmapsto\underbrace{f(x+y)-f(x)-f(y)}_{\varphi_f(x,y)}\}
\]
% archive sheet number: 10
\scanpage{12}

est bilinéaire en $x,y$. \deleted{\unsure{S}}

Comme on a, par transposition de $i:\Delta\longrightarrow\underline{\mathrm{Quad}}(E)$

\[\underline{\mathrm{Quad}}(E)^\vee\xrightarrow{i^\vee}\Delta^\vee,\]

on déduit une forme quadratique composée

\diagram{\begin{tikzcd}[ampersand replacement=\&]
\&\underline{\mathrm{Quad}}(E)^\vee\arrow[dr,"i^\vee"]\&\\
Q_i\ \text{ou }Q_\Sigma:\quad E\arrow[ur,"Q"]\arrow[rr]\&\&\Delta^\vee.
\end{tikzcd}}

Dans le cas d'une structure de similitudes fidèle \inserted{$\Sigma$}, $(\Delta,i)$ ne dépend, à iso. \inserted{can.} unique près, que de $\Sigma$, et il en est donc aussi de $=$ de $Q_i$, ce qui justifie alors la notation $Q_\Sigma$. \deleted{Quand on}

\uline{NB} Quand on se donne une base $e$ de $\Delta$, d'où $\Delta\simeq\Delta^\vee\simeq\underline{\mathcal O}$, alors $Q_\Sigma$ s'identifie à une forme quad. ordinaire sur $E$ i.e. à val. dans $\underline{\mathcal O}$, et ce n'est autre que \deleted{$\unsure{f(x)}$} \inserted{$f=i(e)$.}

% margin note written sideways
\margin{1.5. La donnée d'une structure de similitudes \inserted{(\uline{strictement})} \uline{fidèle} \inserted{sur $E$} \unsure{équivaut|revient} à celle d'un Module inv. $\Delta$ sur $X$, et d'une forme quadratique $Q:E\longrightarrow\Delta^\vee$ telle que \uline{$Q(E)$ engendre} $\Delta^\vee$, i.e. \deleted{\unsure{$Q_x$ est non|$Q$}} sur aucune fibre résiduelle, \deleted{\unsure{est nul}} identiquement nulle.}

2.1. On suppose à partir de maintenant que $E$ est loc. libre de rang 2. Soit

\[L=\underline{\det}E=\bigwedge^2 E\]

son Module déterminantiel, qui est un Module inversible. Notons l'iso. can.

\[E\xrightarrow{\sim} E^\vee\otimes L\]

donné par

\[(x\longmapsto(y\longmapsto x\wedge y)),\quad\text{i.e.}\]

associé à l'accouplement canonique

\[E\times E\longrightarrow L=\bigwedge^2 E:(x,y)\longmapsto x\wedge y.\]
% archive sheet number: 11
\scanpage{13}

On en déduit des iso. can.

\[
\left\{\begin{aligned}
E^\vee&\simeq E\otimes L^\vee\\
\underbrace{\mathrm{Sym}^2(E^\vee)}_{\underline{\mathrm{Quad}}(E)}
&\simeq\underbrace{\mathrm{Sym}^2(E)}_{\underline{\mathrm{Quad}}(E^\vee)}\otimes
\underbrace{L^{\vee\otimes2}}_{L^{\otimes(-2)}}
\end{aligned}\right.
\]

\margin{i.e. $\underline{\mathrm{Quad}}(E)\simeq\mathrm{Quad}(E^\vee)\otimes L^{\otimes\unsure{-2}}$}

Il en résulte un isom. canonique de faisceaux

\[\underline{\mathrm{Sim}}(E)\xrightarrow{\substack{\mathrm{can}_E\\\sim}}\underline{\mathrm{Sim}}(E^\vee)\]

qui transforme le sous-Module loc. monogène $\Delta$ de $\underline{\mathrm{Sym}}^2(E^\vee)\simeq\mathrm{Quad}(E)$ dans

\[
\begin{aligned}
\Delta\otimes L^{\otimes(\unsure{2|-2})}&\subset
\deleted{\underline{\mathrm{Sym}}^2(E)}\ \underline{\mathrm{Quad}}(E^{\deleted{\unsure{\vee}}})\otimes L^{\otimes2}
\simeq \mathrm{Quad}(E^\vee)\otimes L^{\otimes(-2)}\inserted{\unsure{\otimes L^{\otimes2}|\otimes2}}\\
&\simeq\mathrm{Quad}(E^\vee).
\end{aligned}
\]

On trouve que $\mathrm{can}_E$ et $\mathrm{can}_{E^\vee}$ sont inverses l'un de l'autre. Les notions de \inserted{structures de} similitudes \uline{fidèles}, strictement fidèles, non-dégénérées (transversalement, universellement \ldots) \inserted{pour $E$ et pour $E^\vee$} se correspondent.

2.2. Considérons encore l'iso. can. (dépendant du choix d'un signe)

\[E^\vee\simeq E\otimes L^\vee,\]

en \inserted{$\otimes$-}multipliant \inserted{à droite} par $E$, on trouve

\[E^\vee\otimes E\simeq E\otimes E\otimes L^\vee\]

D'autre part on a \deleted{\unsure{un is. canonique}} \inserted{(pour tout $\underline{\mathcal O}_X$-Module)} une suite exacte canonique (qu'on peut mettre un $0$ à g si $E$ loc. Libre \ldots)

\[
\begin{aligned}
\bigwedge^2 E&\longrightarrow E\otimes E\longrightarrow\mathrm{Sym}^2(E)\longrightarrow0,\\
x\wedge y&\longrightarrow x\otimes y-y\otimes x
\end{aligned}
\]
% archive sheet number: 12
\scanpage{14}

et dans le cas actuel $\bigwedge^2 E\overset{\unsure{\mathrm{déf}|\mathrm{dfn}|\mathrm{df}}}=L$ est inversible, \deleted{et de plus} donc

\[0\mathrel{-}L\longrightarrow E\otimes E\mathrel{-}\mathrm{Sym}^2(E)\mathrel{-}0,\]

d'où en multipliant par $L^\vee$

\[0\longrightarrow\underline{\mathcal O}\longrightarrow\underbrace{(E\otimes E)\otimes L^\vee}_{\rotatebox{90}{$\simeq$}\ E^\vee\otimes E\simeq\underline{\mathrm{End}}(E)}\mathrel{-}\mathrm{Sym}^2(E)\otimes L^\vee\mathrel{-}0\]

Je dis que par l'iso. can. vertical, le sous-Module $L\otimes L^\vee\simeq\underline{\mathcal O}$ est transformé en le sous-Module $\underline{\mathcal O}.1_E$ de $\underline{\mathrm{End}}(E)$, de façon plus précise qu'on a \inserted{un diagramme} commutatif \deleted{de \unsure{du|des}} d'iso. canonique

\diagram{\begin{tikzcd}[ampersand replacement=\&,column sep=small,row sep=large]
0\arrow[r,dash]\&\underset{\underline{\mathcal O}}{\underbrace{\overbrace{L}^{\det E}\otimes L^\vee}}\arrow[r]\arrow[d,"\sim"'," {\mathrm{id}_{\underline{\mathcal O}}}"]\&(E\otimes E)\otimes L^\vee\arrow[r,dash]\arrow[d,"\sim"']\&\overbrace{\mathrm{Sym}^2(E)}^{\mathrm{Quad}(E^\vee)\simeq\mathrm{Quad}(E)\otimes L^{\otimes2}}\otimes L^\vee\arrow[r]\arrow[d,"\sim"']\&0\\
0\arrow[r,dash]\&\underline{\mathcal O}\arrow[r,hook]\&\underset{\substack{\simeq\\\underbrace{E^\vee}_{E\otimes L^\vee}\otimes E}}{\underline{\mathrm{End}}(E)}\arrow[r,dash]\&\underline{\mathrm{End}}(E)/\underline{\mathcal O}\arrow[r,dash]\&0,
\end{tikzcd}}

d'où je retiens pour l'instant surtout

\diagram{\begin{tikzcd}[ampersand replacement=\&]
\mathrm{Sym}^2(E)\otimes L^\vee\simeq\mathrm{Quad}(E)\otimes L\quad(\simeq\mathrm{Quad}(E,L))\arrow[d,"\sim"']\\
\underline{\mathrm{End}}(E)/\mathcal O
\end{tikzcd}}
% archive sheet number: 13
\scanpage{15}

On en conclut : \uline{les structures de similitude} sur le Module (loc. libre de rang 2) $E$ correspondent can. aux sous-Modules loc. monogènes du \inserted{Module (loc. libre de rang 3)} $\underline{\mathrm{End}}(E)/\mathcal O$, i.e. aux sous-Modules $\mathcal A$ de $\underline{\mathrm{End}}(E)$ tels que

\[
\begin{aligned}
1^\circ)&\quad\mathcal A\supset\underline{\mathcal O}\cdot1_E\\
2^\circ)&\quad\mathcal A/\underline{\mathcal O}/1_E \text{ soit loc. monogène.}
\end{aligned}
\]

Les structures de similitudes strictement fidèles correspondent aux sous-fibrés vectoriels de rang 2 (du \inserted{fibré vect. de rang 4} $\underline{\mathrm{End}}(E)$), qui contiennent le sous-fibré $\underline{\mathcal O}\cdot1_E$.

\uline{Lemme 1} Soit $\mathcal A$ un sous-Module de $\underline{\mathrm{End}}(E)$ qui contient $\underline{\mathcal O}\cdot1_E$. Conditions équivalentes

a) $\mathcal A/\underline{\mathcal O}1_E$ est monogène

a$'$) $\exists$ loc. une section $u$ de $\mathcal A$ t.q. $\mathcal A=\underline{\mathcal O}\cdot1_E+\underline{\mathcal O}\cdot u$.

b) $\mathcal A$ est engendré loc. par deux sections.

% slanted
\margin{résulte de l'équation de H.C. $u^2-\mathrm{Tr}(u)\,u+N(u)=0$}

Un tel $\mathcal A$ est stable par \deleted{\unsure{la}} la loi d'algèbre de $\underline{\mathrm{End}}(E)$, et est un faisceau d'anneaux comm. Ainsi \inserted{(op. ou plutôt, \unsure{A de plus})} les $A$ en question \deleted{ainsi} font une \unsure{O|0}). \unsure{revient à dire que}

correspondent : des \inserted{structures de} similitudes sur $E$) sont des sous-algèbres \inserted{commutatives} de $\underline{\mathrm{End}}(E)$. Les str. de similitude str. fidèles correspondent loc\textsuperscript{t} aux $\deleted{\mathcal A}\mathcal A_u$ \deleted{loc. définis par $u$} tels que $u$ ne soit
% archive sheet number: 14
\scanpage{16}

un scalaire sur aucune fibre résiduelle.

$A$ est appelé l'anneau des \uline{homothéties} (\inserted{sous-entendu :} \uline{directes}) \deleted{\unsure{de|pour}} pour la structure de similitude envisagée $\Sigma$, on le note $\mathcal A_\Sigma$.

\uline{Lemme 2} Soit $\Sigma$ une structure de similitude sur $E$, et $\mathcal A_\Sigma$ \inserted{$\subset\underline{\mathrm{End}}(E)$} l'Anneau des homothéties. Conditions équivalentes

a) $\deleted{\unsure{\Sigma|\Sigma_u}}\Sigma$ est strictement fidèle.

b) \inserted{$\mathcal A(s)$ est de rang 2 et} $\mathcal A(\deleted{\unsure{x}}s)-\underline{\mathrm{End}}(E(s))$ injectif pour tout $s\in X$, i.e. $\mathcal A$ sous-fibré vectoriel (de rang 2) de $\underline{\mathrm{End}}(E)$.

b$'$) (Si $\mathcal A=\mathcal A_u\overset{\mathrm{déf}}{=}\underline{\mathcal O}\cdot1_E+\underline{\mathcal O}\cdot u$) $\forall s\in X$, $u(s)\in\underline{\mathrm{End}}(E(s))$ n'est pas un scalaire

c) $E$ est un $\mathcal A$-Module \deleted{loc. libre de \unsure{rg|rang}} \inserted{inversible}, i.e loc. libre de rang 1.

c$'$) $E$ est un $\mathcal A$-module \inserted{loc. \deleted{\unsure{mon}}} monogène.

d) $\mathcal E=\underline{\mathrm{End}}(E)$ est un $\mathcal A$-Module \unsure{à|;} g. (pour la structure régulière gauche) loc libre de rang 2.

d$'$) $\mathcal E$ est un $\mathcal A$-module \unsure{loc\textsuperscript{t}|loc ?|loc.} engendré par 2 élém\textsuperscript{ts}.

\uline{Cor. 1}, Sous ces conditions, soit \deleted{$x\in\Gamma E$, alors $x$ est} \inserted{$\Delta\subset\underline{\mathrm{Quad}}(E)$} le sous-Module \deleted{loc. lib.} inv. de $\underline{\mathrm{Quad}}(E)$ tel que $\Sigma=\Sigma_\Delta$. On aura donc un iso. can.

\[
\begin{aligned}
&\Delta\simeq(\mathcal A/\mathcal O)\otimes L^\vee &&(\text{où } L=\Lambda^{\unsure{2|7}}E=\det E)\\
\text{d'où}\qquad&\Delta^\vee\simeq(\mathcal A/\mathcal O)^\vee\otimes L&&,
\end{aligned}
\]
% archive sheet number: 15
\scanpage{17}

et une \deleted{\unsure{ho.|la}} forme quadratique de définition de la structure \deleted{quadratique} \inserted{de} similitude

\[
Q_\Sigma:E\longrightarrow\Delta^\vee\overset{\mathrm{déf}}{=}\Delta' \;(\simeq\mathcal A/\mathcal O\otimes L^\vee)
\]

\deleted{Ceci posé, soit} Plus généralement, soit $\Delta'$ un Module inversible,

\[
Q:E\longrightarrow\Delta'\qquad(\text{i.e. }\Delta^{\prime\vee}\longrightarrow\underline{\mathrm{Quad}}(E)\quad\text{un hom. de Modules})
\]

une forme quadratique, d'où une structure de similitudes $\Sigma=\Sigma_Q$ sur $E$, \unsure{correspond.|correspond\textsuperscript{t}|correspondant} à un sous-anneau $\mathcal A\unsure{=|\supset}\mathcal A_\Sigma$ de $\underline{\mathrm{End}}(E)$.

Soit $x$ une section de $E$. Conditions équivalentes

a) $\{x\}$ est un générateur de $E$ comme $\mathcal A$-Module, i.e. \inserted{$\lambda\mapsto\lambda\cdot x$,} $\mathcal A\longrightarrow\deleted{\unsure{\mathcal A}}E$, est épi

a$'$) \deleted{$Q(x)$ \unsure{es|est}} $\{x\}$ est une base de $E$ comme $\mathcal A$ Module, i.e. \inserted{$u\mapsto u\cdot x$:} $\mathcal A\longrightarrow E$ est un iso

b) $\{Q(x)\}$ est une base de $\deleted{\unsure{E|L}}\Delta'$

b$'$) (Si $\Delta'=\underline{\mathcal O}$) $Q(\unsure{u|x})\in\Gamma\underline{\mathcal O}_X$ est \uline{inversible}.

c) $\forall s\in X$, $Q(s)\ne0$

\deleted{\uline{Cor. 2} Soit $Q$}

d) (Si $u$ est un gén. loc. de l'Alg. $\mathcal A$) Pour \deleted{tout} $s\in X$, $x(s)\in E(s)$ est \uline{non vect.} \uline{propre} de $u(s)\in\mathrm{End}\,E(s)$
% archive sheet number: 16
\scanpage{18}

2.3. Considérons \unsure{à nouveau|: à nouveau} l'iso. can.

\[
\underline{\mathrm{End}}(E)/\underline{\mathcal O}\cdot1_E\simeq\underline{\mathrm{Quad}}(E)\otimes\underline L\;(\simeq\mathrm{Quad}(E,L))
\]

Soit $u$ une section de $\underline{\mathrm{End}}(E)$ i.e. un endom. de $E$, \deleted{considér\unsure{ons}} son image est donc une forme quadratique $Q_u$ à valeurs dans $L$. Je dis qu'on a

\[
\boxed{Q_u(x)=u(x)\wedge x}
\]

(\deleted{mais} il faut vérifier le signe !). \deleted{C'est} On revérifie directement par des coordonnées que

\[
Q_u=0\quad\Longleftrightarrow\quad u\in\Gamma\underline{\mathcal O}\cdot\mathrm{id}_E
\]

[ NB. Si

\[
u=\begin{pmatrix}a&b\\c&d\end{pmatrix}
\]

alors

\[
Q_u(\lambda,\mu)=c\lambda^2+(d-a)\lambda\mu\deleted{\unsure{+|-}}\unsure{-|+}b\mu^2
\]

\deleted{qui} est id. nul en $\lambda,\mu$ ssi $b=c=0$, $d=a$ i.e. $u=\deleted{\unsure{a}}\begin{pmatrix}a&\deleted{\unsure{c|a}}0\\0&a\end{pmatrix}\deleted{\unsure{=a}}=a\,\mathrm{id}_E$

Si $Q=\alpha\lambda^{\unsure{2|7}}+\beta\lambda\mu+\gamma\mu^2$, alors les $u$ tels que $Q_u=Q$ sont ceux de la forme

\[
u_a=\begin{pmatrix}a&-\gamma\\\alpha&a+\beta\end{pmatrix},\quad\text{avec }a\in\Gamma\underline{\mathcal O}\text{ arbitraire}
\]

parmi lesquels

\[
u_0=\begin{pmatrix}0&-\gamma\\\alpha&\beta\end{pmatrix}\qquad\ldots\qquad]
\]

On a \deleted{alors}
% archive sheet number: 17
